% Part: applied-modal-logics
% Chapter: epistemic-logic
% Section: bisimulations

\documentclass[../../../include/open-logic-section]{subfiles}

\begin{document}

\olfileid{aml}{el}{bsd}

\olsection{بائی سیمولیشناں}

ساڈی علم دی زبان دی اظہاری طاقت بارے اک ہور سوال نمونیاں تے
اوہناں وچ لاگو ہون والے فارمولیاں دے تعلق بارے اے۔ فریم خاصیت تے
فارمولے دے میل والے نتیجیاں توں اسیں ویکھیا اے کہ جدوں کجھ فارمولے
کسے فریم وچ معتبر ہون، تاں اوہ ایہ وی یقینی بناؤندے نیں کہ فریم
کجھ خاصیتاں پوریاں کردا اے۔ پر مثلاً کی ساڈی امکانی زبان سانوں
اک ایسے جہان، جیہدے توں اوہدے اپنے ول تیر اے، تے جہاناں دی اک
لامتناہی لڑی، جیہدے وچ ہر جہان توں اگلے ول تیر اے، وچ فرق دسن
دیندی اے؟ یعنی کی کوئی ایسا فارمولا $A$ اے جیہڑا ایہناں دو جہاناں
وچوں صرف اک اُتے لاگو ہو سکدا اے؟

بائی سیمولیشن اک تعلق اے جیہڑا تعلقی نمونیاں وچ ایہ دسن لئی
متعین کر سکدے آں کہ اوہناں دی بناوٹ اثر دے لحاظ نال اکو جیہی اے۔
اسیں ویکھاں گے کہ ایہ نمونیاں دی اوہ ہم ارزی پھڑدا اے جیہڑی ساڈی
علم دی زبان وچ ظاہر ہو سکدی اے۔

\begin{defn}[بائی سیمولیشن]
منّ لو $M_1 = \tuple{W_1, R_1, V_1}$ تے $M_2 = \tuple{W_2, R_2, V_2}$
دو تعلقی نمونے نیں۔ منّ لو $\mathcal{R} \subseteq W_1 \times W_2$
اک دو رکنی تعلق اے۔ $\mathcal{R}$ نوں \emph{بائی سیمولیشن} آکھدے
آں جدوں ہر $\tuple{w_1, w_2} \in \mathcal{R}$ لئی ایہ ہووے:

\begin{enumerate}
  \item $w_1 \in V_1(p)$ اُدوں تے صرف اُدوں جدوں $w_2 \in V_2(p)$،
    سارے !!{propositional variable}s~$p$ لئی۔
  \item سارے فاعلاں $a \in A$ تے جہاناں $v_1 \in W_1$ لئی، جے
    $R_{1_a} w_1 v_1$، تاں کوئی $v_2 \in W_2$ اے جیہدے لئی
    $R_{2_a} w_2 v_2$ تے $\tuple{v_1, v_2} \in \mathcal{R}$۔
  \item سارے فاعلاں $a \in A$ تے جہاناں $v_2 \in W_2$ لئی، جے
    $R_{2_a} w_2 v_2$، تاں کوئی $v_1 \in W_1$ اے جیہدے لئی
    $R_{1_a} w_1 v_1$ تے $\tuple{v_1, v_2} \in \mathcal{R}$۔
\end{enumerate}

جدوں $M_1$ تے $M_2$ وچ کوئی بائی سیمولیشن اے جیہڑی جہاناں
$w_1$ تے $w_2$ نوں جوڑدی اے، تاں اسیں
$\tuple{M_1, w_1} \leftrightarroweq \tuple{M_2, w_2}$ وی لکھ سکدے
آں، تے $\tuple{M_1, w_1}$ تے $\tuple{M_2, w_2}$ نوں
\emph{بائی سیمولیشن نال جڑے} آکھدے آں۔
\end{defn}

بائی سیمولیشن تعلق دیاں وکھ وکھ شقّاں وکھ وکھ گلاں یقینی بناؤندیاں
نیں۔ شق 1 یقینی بناؤندی اے کہ بائی سیمولیشن نال جڑے جہان امکانی
عامل توں خالی اوہی فارمولے پورے کرن گے، کیوں جے سارے
!!{propositional variable}s اُتے میل اے۔ دوجیاں دو شقّاں نوں کدی
کدی واری واری «اگے» تے «پچھے» آکھدے نیں۔ اوہ یقینی بناؤندیاں نیں
کہ رسائی دے تعلقاں دی بناوٹ اکو جیہی ہووے گی۔

\begin{thm}
جے $\tuple{M_1, w_1} \leftrightarroweq \tuple{M_2, w_2}$، تاں ہر
!!{formula}~$!A$ لئی $\mSat{M_1}{!A}[w_1]$ اُدوں تے صرف اُدوں
جدوں $\mSat{M_2}{!A}[w_2]$۔
\end{thm}

\begin{figure}
  \begin{center}
    \begin{tikzpicture}[modal]
      \node[world] (w1) {$w_1$}; 
      \node[world] (w2) [above left=of w1]{$w_2$}; 
      \node[world] (w3) [above right=of w1] {$w_3$};
      \draw[<->] (w1) to node {$a$} (w2);
      \draw[->,loop below] (w1) to node {$a$} (w1);
      \draw[<->] (w1) to [swap] node {$a$} (w3);
      \draw[->,loop above] (w2) to node {$a$} (w2) ;
      \draw[->,loop above] (w3) to node {$a$} (w3) ;
      
      \node[world](v1)[right of=w1, xshift=1.5in]{$v_1$};
      \node[world](v2)[above of=v1]{$v_2$};
      \draw[<->] (v1) to node {$a$} (v2);
      \draw[->,loop below] (v1) to node {$a$} (v1);
      \draw[->,loop above] (v2) to node {$a$} (v2) ;

      \draw[-,dotted] (w1) to (v1) ;
      \draw[-,dotted,bend left=45] (w2) to (v2) ;
      \draw[-,dotted,bend left=30] (w3) to (v2) ;
    \end{tikzpicture}
  \end{center}
  \caption{بائی سیمولیشن نال جڑے دو نمونے۔}
  \ollabel{fig:bisimilar}
\end{figure}

\olref{fig:bisimilar} وچ وکھائے دونیں نمونے کجھ وکھرے نیں، پر
اک بائی سیمولیشن جہاناں $w_1$ تے~$v_1$ نوں جوڑدی اے۔ ایہ بائی
سیمولیشن $w_2$ تے $w_3$ دونہاں نوں~$v_2$ نال وی جوڑے گی۔
گل ایہ اے کہ ساڈی امکانی زبان وچ لکھن جوگ کوئی شے ایہناں وچ فرق
نہیں دس سکدی۔ پر جے $w_2$ تے~$w_3$ وکھ وکھ
!!{propositional variable}s پورے کردے ہون، تاں صورت وکھری ہووے گی۔

\end{document}
